\documentclass[11pt,a4paper]{article}

\usepackage[T1]{fontenc}
\usepackage[utf8]{inputenc}
\usepackage{lmodern}
\usepackage[margin=2.4cm]{geometry}
\usepackage[hidelinks]{hyperref}
\usepackage{booktabs}
\usepackage{tabularx}
\usepackage{enumitem}
\usepackage{parskip}
\usepackage{array}

\setlength{\parindent}{0pt}
\setlist[itemize]{leftmargin=1.5em}
\setlist[enumerate]{leftmargin=1.7em}
\newcolumntype{Y}{>{\raggedright\arraybackslash}X}

\title{\textbf{ReconcileUnderPolicy: A Pilot Benchmark for Policy-Constrained Recovery in the RabbitMQ Cluster Operator}}
\author{Mandar Joshi \\ \texttt{mandarjoshi1045@gmail.com}}
\date{April 2026}

\begin{document}

\maketitle

\begin{abstract}
Kubernetes operators encode operational knowledge into control loops that
deploy, scale, update, and recover complex applications. Prior work has shown
that operator bugs are common and that controller correctness is challenging to
validate, but fewer studies provide concrete, reproducible measurements of how a
real operator behaves under policy-constrained recovery scenarios. This paper
presents \emph{ReconcileUnderPolicy}, a small but executable benchmark focused on
stateful operator recovery. The benchmark is implemented in a standalone
repository, vendors the RabbitMQ Cluster Operator manifests locally, and runs on
a single-node \texttt{kind} cluster. We evaluate three scenarios against the
official RabbitMQ Cluster Operator: baseline cluster creation, scale-up during
operator restart, and scale-up blocked by a temporary ResourceQuota. All three
scenarios completed successfully in the pilot experiment, with end-to-end
recovery times of 64.83\,s, 79.82\,s, and 79.49\,s respectively. The results
do not yet constitute a broad operator benchmark across multiple systems, but
they demonstrate that the proposed methodology is implementable, produces real
artifact-backed measurements, and can distinguish ordinary reconciliation from
fault- and policy-aware recovery workflows.
\end{abstract}

\section{Introduction}

Kubernetes operators extend the platform with custom resources and controllers
that automate domain-specific operational tasks such as provisioning,
reconfiguration, scaling, and cleanup~\cite{k8s-operator}. In modern
cloud-native environments, operators are increasingly responsible for stateful
infrastructure such as databases, queues, and messaging systems. This makes
their recovery behavior a systems concern rather than merely an implementation
detail.

The difficulty of building reliable operators is now well established. Recent
work has demonstrated the prevalence of operator bugs, the complexity of
controller interactions, and the value of systematic testing and fault
injection~\cite{acto,sieve,watchers,issta-bugs}. However, many papers in this
area either study broad bug corpora or introduce sophisticated frameworks that
require substantial infrastructure to evaluate. For an individual researcher or
student project, a practical question remains: can we build a small but real
benchmark that exercises policy-constrained recovery on an actual operator and
produces defensible measurements?

This paper answers that question with a focused pilot study. We implement
\emph{ReconcileUnderPolicy}, a benchmark repository that targets the official
RabbitMQ Cluster Operator~\cite{rabbitmq-quickstart}. The benchmark vendors the
operator manifests locally, launches a dedicated \texttt{kind} cluster, and
runs three scenarios:

\begin{enumerate}
  \item baseline cluster creation,
  \item scale-up while deliberately restarting the operator pod, and
  \item scale-up temporarily blocked by a ResourceQuota and resumed after the
  quota is removed.
\end{enumerate}

The benchmark therefore moves beyond a paper-only proposal. It includes:

\begin{itemize}
  \item a runnable Python benchmark harness,
  \item vendored operator manifests,
  \item scenario-specific snapshots of pod, StatefulSet, cluster, and event
  state,
  \item a machine-readable summary of measured timings and final outcomes.
\end{itemize}

This is still a pilot rather than a full multi-operator study. The present
artifact covers a single operator and only one trial per scenario. Even so, it
demonstrates a viable end-to-end methodology for policy-constrained recovery
evaluation and provides a concrete base for a more ambitious empirical paper.

\section{Background and Motivation}

Operators follow the Kubernetes reconciliation model: a controller observes the
current cluster state, compares it to the desired state expressed in a custom
resource, and issues actions to reduce any mismatch~\cite{k8s-operator}. For
stateful systems, this process involves more than creating Pods. It often
includes storage provisioning, identity management, service configuration, and
ordered startup or shutdown sequences.

Failures during recovery are particularly important. An operator may be correct
for a clean create or update path yet fail when:

\begin{itemize}
  \item the controller itself restarts mid-reconcile,
  \item a cluster policy prevents creation of the next required object,
  \item a stateful workload is left in a partial intermediate state,
  \item cleanup and restart order matter.
\end{itemize}

The pilot benchmark focuses on two of these recovery axes:

\begin{itemize}
  \item \textbf{controller interruption}: deleting the operator pod during a
  scale-up event;
  \item \textbf{policy-constrained recovery}: using ResourceQuota to block a
  scale-up and then removing the constraint to observe convergence.
\end{itemize}

\section{Benchmark Design}

\subsection{Repository Layout}

The benchmark is implemented as a standalone repository to avoid coupling the
artifact to unrelated project code. The relevant components are:

\begin{itemize}
  \item \texttt{vendor/cluster-operator-main}: locally vendored RabbitMQ Cluster
  Operator source tree;
  \item benchmark runner script for the RabbitMQ scenarios;
  \item \texttt{benchmark/reconcileunderpolicy/results/}: generated manifests,
  scenario snapshots, and summary JSON.
\end{itemize}

\subsection{Scenarios}

The pilot study runs the following three scenarios:

\begin{enumerate}
  \item \textbf{baseline\_create}: create a single-replica RabbitMQ cluster and
  measure time to first healthy reconciliation;
  \item \textbf{restart\_during\_scale}: scale from one to two replicas while
  deleting the operator pod during reconciliation;
  \item \textbf{quota\_blocked\_scale}: enforce a pod-count ResourceQuota to
  block creation of the third pod during scale-up, then remove the quota and
  measure time to converge to three healthy replicas.
\end{enumerate}

\subsection{Execution Environment}

The experiment was run on a dedicated \texttt{kind} cluster named
\texttt{reconcile-under-policy}. The cluster used Kubernetes v1.34.0 and the
RabbitMQ Cluster Operator install was built locally with
\texttt{kubectl kustomize} from the vendored operator repository. The managed
RabbitMQ image was \texttt{rabbitmq:4.1.3-management}, and the benchmark
overrode the workload resources to use:

\begin{itemize}
  \item CPU request: 200m
  \item Memory request: 512Mi
  \item CPU limit: 500m
  \item Memory limit: 1Gi
\end{itemize}

These lower values were necessary because the operator defaults are too heavy
for a single-node pilot cluster with three replicas.

\subsection{Measurements}

Each scenario records:

\begin{itemize}
  \item end-to-end completion time in seconds,
  \item final number of ready replicas,
  \item scenario-specific metadata, such as the deleted operator pod name or
  the number of ready replicas observed while the quota constraint was active.
\end{itemize}

The current pilot uses one run per scenario. Therefore, the measurements should
be interpreted as proof-of-execution and initial signals, not as statistically
stable performance claims.

\section{Results}

Table~\ref{tab:results} summarizes the measured outcomes.

\begin{table}[ht]
\centering
\caption{Pilot benchmark results for the RabbitMQ Cluster Operator}
\label{tab:results}
\begin{tabularx}{\linewidth}{@{}lcccY@{}}
\toprule
\textbf{Scenario} & \textbf{Success} & \textbf{Time (s)} & \textbf{Final Ready} & \textbf{Observation} \\
\midrule
baseline\_create & Yes & 64.83 & 1 & First healthy reconciliation from empty state \\
restart\_during\_scale & Yes & 79.82 & 2 & Operator pod deleted during scale-up; cluster still converged \\
quota\_blocked\_scale & Yes & 79.49 & 3 & Ready replicas stayed at 2 while quota blocked progress, then converged after quota removal \\
\bottomrule
\end{tabularx}
\end{table}

Three observations stand out.

\textbf{Baseline startup dominates the first scenario.} The baseline time of
64.83\,s includes image startup, the operator-driven setup sequence, and the
RabbitMQ init-container delay used by the operator-managed workload.

\textbf{Operator restart did not prevent convergence.} In the second scenario,
the benchmark deleted the operator pod during scale-up from one to two replicas.
The cluster still reached the desired healthy state, indicating that the RabbitMQ
operator tolerated this interruption for the tested workflow.

\textbf{Policy-constrained recovery was observable and reversible.} In the third
scenario, the benchmark applied a ResourceQuota that prevented creation of the
third pod during scale-up to three replicas. While the quota was active, the
cluster remained at two ready replicas. After the quota was removed, the
operator progressed to the intended three-replica healthy state. This is
exactly the class of behavior the benchmark is meant to expose.

\section{Discussion}

\subsection{What This Pilot Demonstrates}

The main contribution of the pilot is methodological. It shows that a
policy-constrained recovery benchmark for a real stateful operator can be:

\begin{itemize}
  \item implemented as a standalone repository,
  \item executed end-to-end on a local \texttt{kind} cluster,
  \item backed by machine-readable results and scenario snapshots,
  \item used to observe both controller-interruption and policy-blocked recovery.
\end{itemize}

This is stronger than a purely conceptual proposal because the artifact now
exists and produces real outputs. The benchmark also surfaced practical lessons:

\begin{itemize}
  \item local clusters require careful resource tuning;
  \item preloading images materially improves repeatability;
  \item policy constraints such as ResourceQuota are a useful way to create
  realistic, operator-visible recovery stalls without modifying operator code.
\end{itemize}

\subsection{Limitations}

This pilot is not yet a full publication-ready empirical study for a major
systems venue.

The most important limitations are:

\begin{itemize}
  \item only one operator was evaluated;
  \item only three scenarios were executed;
  \item each scenario was run once, so timing variance is unknown;
  \item the cluster was a single-node \texttt{kind} environment rather than a
  multi-node or production-like deployment.
\end{itemize}

Accordingly, the present results are best positioned as a pilot study, artifact
demo, poster, or workshop paper foundation rather than a finished top-tier
conference submission.

\subsection{Next Steps}

The benchmark can be extended in several straightforward ways:

\begin{enumerate}
  \item add repeated trials per scenario;
  \item add deletion/cleanup recovery scenarios;
  \item introduce a second operator, such as PostgreSQL or Redis;
  \item generate comparative tables and simple charts from the JSON results;
  \item evaluate more policy mechanisms, such as RBAC restriction or admission
  rejection.
\end{enumerate}

\section{Related Work}

Acto~\cite{acto} introduced automatic end-to-end testing for Kubernetes
operators and demonstrated that state-centric testing can uncover many operator
correctness and recovery bugs. Sieve~\cite{sieve} explored perturbation-based
controller testing and emphasized the importance of intermediate, stale, and
unobserved states. The NSDI 2026 reliability study~\cite{watchers} and the
ISSTA 2024 operator bug study~\cite{issta-bugs} provide broader empirical
motivation by showing that operator failures are common and operationally
significant.

The present work is much smaller in scope than those systems papers. Its value
lies in making the policy-constrained recovery idea concrete with a runnable
artifact and actual measured scenarios in a standalone repository.

\section{Conclusion}

This paper presented \emph{ReconcileUnderPolicy}, a pilot benchmark for
policy-constrained recovery in the RabbitMQ Cluster Operator. The benchmark was
implemented, executed, and produced real results across three scenarios:
baseline creation, operator-restart scale-up, and ResourceQuota-constrained
recovery. All three scenarios converged successfully, with measured end-to-end
times between 64.83\,s and 79.82\,s.

The pilot does not yet establish a broad empirical characterization of
stateful Kubernetes operator recovery. What it does establish is that the
benchmarking approach is feasible, reproducible, and capable of observing
policy-aware recovery behavior on a real operator. That makes it a practical
foundation for a more complete multi-operator study.

\bibliographystyle{plain}
\bibliography{references}

\end{document}
